\subsection*{MT-M24. The specified Masullo smooth quotient is the Meyer quotient}
This is a comparison of the actual test spaces and summation maps. It does
not assume Masullo's global converse theorem, his spectral basis claims,
or $M$. Fix the primitive unitary Hecke function and the convention
$\beta(\mathfrak p_v)^{-1}=c_{\mathfrak p_v}$ from MT-M22; put
$C=C_K$, $C^1=\ker|\cdot|$, and $\chi=\beta|_{C^1}$.
Use the norm splitting and volume-one measure on $C^1$ chosen there.
Write $\mathcal E$ for the exponential Schwartz space of MT-M18,
$\mathcal I_\eta=\overline{\psi_\eta*\mathcal E}$ for its closed
completed-Hecke ideal, and $Q_\eta=\mathcal E/\mathcal I_\eta$.
For a real number $a$, let $M_a u(g)=|g|^a u(g)$, and define
\[
 \Theta_\chi=\Phi_{K,\chi}\bigl(\mathcal S_0(\mathbf A_K)\bigr),
 \qquad
 Q^{\mathrm{Mas}}_{\chi,a}
 =\mathcal S_{\mathrm E}(C)(\chi)\big/
       \overline{M_a\Theta_\chi}^{\,\mathcal S_{\mathrm E}(C)(\chi)}.
 \tag{M24.1}
\]
Here the symbols on the right have their literal meanings in Masullo:
\[
 \mathcal S_0(\mathbf A_K)
  =\{f\in\mathcal S(\mathbf A_K):f(0)=\widehat f(0)=0\},
 \qquad
 \Phi_K(f)(g)=\sum_{\xi\in K^\times}f(\xi g),
 \tag{M24.2}
\]
and $\Phi_{K,\chi}$ is compact-character averaging of $\Phi_K$.
The closure in (M24.1) does not enlarge $M_a\Theta_\chi$.
There is an explicit topological isomorphism
\[
 Q_\eta\xrightarrow{\ \sim\ }Q^{\mathrm{Mas}}_{\chi,a},
 \qquad
 [f]\longmapsto
 \bigl[g\longmapsto |g|^a\beta(g)^{-1}f(\log|g|)\bigr].
 \tag{M24.3}
\]
For $a=0$ this identifies the quotient in Masullo (7.1.4.7),
original PDF p.~327. For $a=1/2$ it identifies the centered quotient
in (1.8.3.1), original PDF p.~45. These two source displays use
different twists; (M24.3) keeps the twist explicit.

\begin{proof}
First verify the test-space topology rather than just its elements.
Masullo Definition D.3.3.1, original PDF p.~804, requires
$|\cdot|^b u\in\mathcal S(C)$ for every real $b$.
Definition D.3.3.3, p.~805, gives each fixed compact-invariance
level the projective topology of all these weighted Schwartz spaces,
and gives their union the locally convex inductive topology.
On the $\chi$-sector the map
\[
 \mathcal L_\beta f(g)=\beta(g)^{-1}f(\log|g|)
 \tag{M24.4}
\]
is a linear bijection from $\mathcal E$ onto that sector.
Indeed, the $\chi$-eigencondition for left translation says
$u(cg)=\chi(c)^{-1}u(g)$ for $c\in C^1$; its restriction to
the norm splitting therefore determines $u$. The unitary value
of $\beta$ on that splitting is $e^{i\tau x}$ for a fixed real
$\tau$, so removing it only changes derivatives by fixed finite
linear combinations.

Every function in this sector is invariant under $\ker\chi$.
The compact quotient $C^1/\ker\chi$ is a closed subgroup of the
circle, hence is either finite or a circle, and this is one permitted
elementary compact-invariance level. At this level the weighted
Schwartz seminorms, after $x=\log|g|$, are equivalent to
\[
 p_{N,j}(f)=\sup_x e^{N|x|}|f^{(j)}(x)|,
 \qquad N,j\ge0.
 \tag{M24.5}
\]
To see the equivalence directly, the weights $b=N$ and $b=-N$
bound (M24.5); conversely every polynomial factor in $x$ occurring
in a weighted Schwartz seminorm is absorbed by one larger
exponential weight. Derivatives on the compact factor multiply
this fixed character by constants. Inclusion of this level in the
inductive limit is continuous. The inverse of (M24.4) on the
sector is also continuous: compact-character averaging followed
by restriction to the norm splitting is continuous into
$\mathcal E$ on each elementary level, by these same seminorm
estimates. It is therefore continuous on the inductive limit.
This establishes the stated sector topology directly, without
using a topological claim about a different quotient.

Next identify the denominator as a set of actual functions.
Use Meyer's common ambient space and embeddings from MT-M22:
\[
 i_+(v)=(\Sigma v,J\Sigma\mathcal F_*v),
 \qquad i_-(u)=(u,u),
 \qquad H_+=\mathcal S(\mathbf A_K)/K^\times,
 \quad H_-=\mathcal S_{\mathrm E}(C).
\]
The summation maps agree term by term:
$\Sigma f(g)=\Phi_K(f)(g)$ by (M24.2).
Their Poisson identity is
\[
 J\Sigma\widehat f-\Sigma f
      =f(0)-|\cdot|^{-1}\widehat f(0).
 \tag{M24.6}
\]
For $f\in\mathcal S_0(\mathbf A_K)$ the two entries agree.
The first entry has all the right-half-plane Schwartz weights,
and the second all the left-half-plane weights, by the defining
Meyer embeddings. Their common value has all real weights:
between any two admitted weights the weighted derivative is
bounded by the two endpoint weighted derivatives, and polynomial
Schwartz factors are treated in the same way. Thus this common
value belongs to $H_-$, and
$\Phi_K(\mathcal S_0)\subset H_+\cap H_-$.

Conversely, if $(u,u)=i_+(v)$ lies in this intersection, choose
an adelic Schwartz representative $f$ of the coinvariant class
$v$. Equation (M24.6) then vanishes identically. The functions
$1$ and $|\cdot|^{-1}$ are linearly independent because the norm
map is onto $\mathbf R_{>0}$, so $f(0)=\widehat f(0)=0$.
Consequently
\[
       \Phi_K(\mathcal S_0(\mathbf A_K))=H_+\cap H_-.
 \tag{M24.7}
\]
This argument uses an actual representative of the coinvariant
class, not a claim that every completion class has an adelic
representative. $H_+$ here is Meyer's quotient of the adelic
Schwartz space.

Compact-character averaging preserves both embedded spaces
and their intersection. On functions, the volume-one projector is
\[
 P_\chi u(g)=\int_{C^1}u(cg)\chi(c)\,dc.
\]
Masullo (2.2.5.2), p.~95, is precisely this formula applied to
$\Phi_K f$, except that his chosen compact Haar measure can have
volume $c_K\ne1$. This multiplies the operator by the fixed
nonzero scalar $c_K$ and leaves its image unchanged. Hence
\[
       \Theta_\chi=(H_+\cap H_-)^\chi.
 \tag{M24.8}
\]
One can also verify surjectivity in (M24.8) directly: an element
of the intersection in the $\chi$-sector already has an
$\mathcal S_0$ representative by (M24.7), and applying
$\Phi_{K,\chi}$ to that representative returns $c_K$ times the
element. No choice of a compact subgroup lift in the idele group
is required.

MT-M22, whose image criterion and topology were proved using
Meyer's specified maps, now gives
\[
 \mathcal L_\beta^{-1}\Theta_\chi
 =\{f\in\mathcal E:(\mathcal Mf)^{(j)}(\rho)=0
         \text{ for all }0\le j<m_\rho\}
 =\mathcal I_\eta.
 \tag{M24.9}
\]
The middle space is closed, since each indicated evaluation is
continuous. Thus $\Theta_\chi$ is closed in the actual sector
topology just checked. Multiplication by $|\cdot|^a$ is an
automorphism of that topology, with inverse multiplication by
$|\cdot|^{-a}$; this follows either from (M24.5), using an
integer larger than $|a|$, or directly by translating the family
of all real weights. This proves the closure assertion and the
continuous two-sided quotient isomorphism (M24.3).
\end{proof}

The action is also explicit. For $b\in C$,
\[
 \lambda_b M_a\mathcal L_\beta f
  =\beta(b)|b|^{-a}M_a\mathcal L_\beta
       \bigl(f(\,\cdot-\log|b|)\bigr).
 \tag{M24.10}
\]
For an idele-class convolution test $k$, set
\[
 k_\beta(y)=\int_{C^1}k(c\tau(e^y))\beta(c\tau(e^y))\,dc.
\]
Changing variables in the Haar integral gives
\[
 C_k M_a\mathcal L_\beta f
   =M_a\mathcal L_\beta\bigl((e^{-a\cdot}k_\beta)*f\bigr).
 \tag{M24.11}
\]
In particular the spectral multiplier is evaluated at $\rho-a$,
with the compact-character factor $\beta$ retained. At $a=0$ this
is precisely the MT-M22 module comparison; at $a=1/2$ it is the
centered version, not an assertion that the same convolution
operator is unchanged by the twist.
An alternative Haar normalization rescales the scalar convolution
integral by the same fixed positive constant; it changes neither
the denominator nor the quotient topology.

\paragraph{Exact source scope and the remaining completion question.}
The original Masullo PDF pages used here are: the theta-sum
definition (2.2.3.1), pp.~81--82, and its underlying sum (2.2.2.1),
p.~77; the polar hyperplanes (2.2.4.23), p.~91;
character averaging (2.2.5.2), p.~95; the actual smooth quotient
(7.1.4.7), p.~327; its centered introductory version (1.8.3.1),
p.~45; and Definitions D.3.3.1/D.3.3.3, pp.~804--805.
The original p.~327 quotient formula was rendered and read;
the other defining formulas were checked against extraction
from the original PDF. The Meyer inputs are his summation map
and Lemma~5.3, pp.~39--40; the embeddings in Lemma~5.4,
pp.~40--41; and Poisson formula (31), p.~41, with the complete
provided arguments already read for MT-M22. No new proof of
the full Masullo draft is claimed.

Masullo Definition~10.1.1.1, p.~434, instead closes
$M_a\Theta_\chi$ in a weighted Hilbert norm, and
(10.2.1.1)--(10.2.1.2), p.~439, and (10.4.1.3), p.~456,
form the resulting Hilbert quotients. Those are different
topological completions of the identified smooth quotient.
Their kernel, retained jets and trace-class properties must be
proved in that norm; MT-M25 and MT-M26 address those maps.
The quotient identification itself is unconditional. Any
subsequent use of universal Hecke GRH is still the conditional
input $M$.

The geometric crosswalk is now an actual Tate--Poisson
summation quotient, rather than an unspecified space sharing
its zero multiplicities. It advances the inherited cote~1 /
SA23 geo-H-lidskii and geo-H-genus interfaces at the precise
test-space and closed-denominator level. It makes no claim
to construct an $\ell$-adic or motivic realization.
